\section{Memory, continual learning, interpretability y alignment}

El capítulo anterior discutió las representaciones erróneas de los modelos sobre el mundo actual. Este capítulo indaga cómo los sistemas guardan información a lo largo del tiempo, actualizan sus parámetros y mantienen la coherencia de objetivos en nuevos entornos. Memory, continual learning, interpretability y alignment suelen aparecer juntos porque abordan conjuntamente "en qué se convierte un sistema tras operar por periodos prolongados".

\subsection{Long-term memory: no es solo extender el contexto}

La memoria a largo plazo incluye al menos escritura, recuperación, actualización, resolución de conflictos, compresión, olvido y permisos. Puede denominar a store de memoria $M_t$, nueva observación como $o_t$ y estrategia de actualización como
\[
M_{t+1}=U(M_t,o_t;\rho,\kappa),
\]
donde $\rho$ es la estrategia de retención/eliminación y $\kappa$ las reglas de resolución de conflictos. Si un hecho nuevo contradice uno antiguo, el sistema puede sobrescribirlo, conservar versiones, reducir la confianza o solicitar confirmación; guardar "todo" sin una política genera ruido y riesgos de privacidad.

\lecturefigure{slide-139.jpg}{Objetivo de la memoria a largo plazo: almacenamiento, recuperación y personalización entre sesiones}{CS25 V6 oficial deck, p.~139}

\lecturefigure{slide-140.jpg}{Métodos de memoria: bases vectoriales, resúmenes, compresión, actualización y olvido}{CS25 V6 oficial deck, p.~140}

\teachervoice{01:03:52--01:06:32, el docente enfatiza que la dificultad de memory no radica en "si se puede guardar", sino en contradicción, pruning y relevancia. Experiencia práctica: un contexto sin límites no equivale a un gestor de memoria confiable.}

\begin{lstlisting}[language=Python,caption={Pseudocódigo de actualización de memoria con estrategia de conflicto y eliminación}]
def update_memory(memory, observation, policy):
    candidates = retrieve_conflicts(memory, observation)
    resolution = policy.resolve(observation, candidates)
    memory.apply(resolution)
    if memory.size > policy.capacity:
        memory.evict(policy.rank_for_eviction(memory))
    return memory
\end{lstlisting}

\begin{knowledgebox}{Digestión de términos de Memory}
\textbf{Context memory} son los tokens presentes en el prompt actual; \textbf{episodic memory} guarda eventos pasados; \textbf{semantic memory} almacena hechos o conceptos; \textbf{vector store} es la implementación de recuperación; \textbf{parameter memory} representa las regularidades estadísticas escritas en los pesos. Difieren en velocidad de actualización, trazabilidad del origen y mecanismos de olvido.
\end{knowledgebox}

\subsection{Aprendizaje a lo largo de la vida: ¿cuándo se considera que un modelo realmente aprendió?}

El curso adopta una definición estricta: RAG, context memory y distillation pueden mejorar el uso continuo, pero si los parámetros del modelo y su estructura de habilidades no se adaptan en línea, no necesariamente se llama true continual learning. Los humanos pueden actualizarse mediante interacción continua, mientras que los LLM suelen depender de reentrenamiento offline periódico. Las métricas clave incluyen plasticity, retention, forward transfer, backward transfer y catastrophic forgetting.

\lecturefigure{slide-141.jpg}{Aprendizaje a lo largo de la vida como los humanos: objetivos y dificultades tras el despliegue}{CS25 V6 oficial deck, p.~141}

\lecturefigure{slide-142.jpg}{Trabajo actual vs. necesario: prompt memory, distillation, edición de modelos y adaptación continua real}{CS25 V6 oficial deck, p.~142}

\teachervoice{01:06:32--01:08:08, el docente distingue claramente entre memory/RAG y continual learning a nivel de parámetros. Nota de clase: colocar información en una base vectorial puede actualizar la fuente de conocimiento, pero no altera la dinámica de aprendizaje del modelo en sí.}

\subsection{Edición de modelos: ¿por qué los cambios locales de pesos causan efectos colaterales?}

Métodos como Rank-One Model Editing (ROME) intentan modificar hechos específicos mediante actualizaciones de rango bajo. Abstractamente, si los pesos de una capa son $W$, la actualización es
\[
W'=W+uv^{\top},
\]
donde $u,v$ controlan la dirección de escritura. Un rango bajo local no equivale a semántica local: un mismo parámetro participa en múltiples contextos, por lo que modificar una relación puede afectar paráfrasis, hechos conexos o comportamientos irrelevantes. La evaluación requiere eficacia, generalización, especificidad y portabilidad, no solo observar el prompt objetivo.

\lecturefigure{slide-143.jpg}{Edición de modelos para aprendizaje continuo: escritura de hechos local y sus efectos secundarios}{CS25 V6 oficial deck, p.~143}

\begin{warningbox}{La edición de modelos no es una UPDATE de base de datos}
Los pesos son representaciones distribuidas sin "ranuras" de hechos independientes. Aunque se corrija la pregunta objetivo, también podría destruir conceptos adyacentes, ser válida solo para una formulación específica o entrar en conflicto con el entrenamiento posterior. Los sistemas de edición deben registrar provenance, versiones y rollback, no considerar un éxito único como aprendizaje seguro en línea.
\end{warningbox}

Los sistemas a largo plazo también deben definir la prioridad entre memory y actualización de parámetros. La información de alta vigencia, revocable y con trazabilidad es más adecuada para store externos; las regularidades de uso repetido, estables y ampliamente transferibles son las que merecen escribirse en los parámetros. Antes de actualizaciones en línea, se debe validar locality, retention y seguridad en un shadow model, y luego publicar mediante rollout versionado. Si retrieval, memory y pesos guardan versiones distintas de hechos, el sistema debe establecer quiénes tienen autoridad ante conflictos y exponer el origen al usuario. Por ello, la dificultad del aprendizaje continuo no es un solo paso de gradientes, sino el ciclo completo: recopilación, validación, escritura, monitoreo, rollback y olvido.

\subsection{Interpretability: entender los mecanismos internos más allá de las salidas}

Interpretability intenta identificar qué representa el modelo, qué circuitos generan la salida y si los fallos pueden detectarse anticipadamente. Behavioral eval solo observa entradas y salidas; mechanistic interpretability analiza activación, características, atención/circuitos e intervención. La evidencia fuerte proviene del cambio en las salidas tras modificar variables internas según predicción, no solo de visualizaciones correlacionales.

\lecturefigure{slide-145.jpg}{Interpretabilidad de LLM: desde caja negra hasta feature, circuitos y detección de fallos}{CS25 V6 oficial deck, p.~145}

\begin{importantbox}{Gradiente de evidencia explicativa}
Visualizar atención es descripción; probe es evidencia de decodificabilidad; activación patching/ablation es evidencia de intervención; la replicación entre muestras y modelos respalda afirmaciones mecanísticas más fuertes. Los métodos explicativos también pueden omitir cómputo distribuido, redundante o dependiente del contexto.
\end{importantbox}

\subsection{Problema de alignment: brechas entre objetivos, comportamiento y nuevos entornos}

Alignment exige que los sistemas persigan objetivos correctos y cumplan restricciones en distintos entornos. Formalmente, la recompensa de entrenamiento $R_{\mathrm{train}}$ es solo un proxy para la intención humana $U$. Si una política maximiza
\[
\pi^{\star}=\arg\max_{\pi}\mathbb{E}[R_{\mathrm{train}}],
\]
pero $R_{\mathrm{train}}\neq U$, pueden surgir reward hacking, shortcut, deception o distribution shift. El problema no es simplemente que el modelo "no obedece", sino que la representación del objetivo y la cobertura de supervisión son incompletas.

\lecturefigure{slide-146.jpg}{El problema de alignment: diferencia entre objetivos de optimización del modelo y valores humanos}{CS25 V6 oficial deck, p.~146}

\lecturefigure{slide-147.jpg}{Modos de fallo: reward hacking, deception, shortcut y distribution shift}{CS25 V6 official deck, p.~147}

\teachervoice{01:08:58--01:11:00, el docente divide el riesgo de alignment en shortcut, reward hacking, deception, post-hoc explanation y desajuste en nuevos entornos. Nota de clase: "comportamiento correcto" en el conjunto de entrenamiento no implica que el objetivo se haya interiorizado correctamente.}

\subsection{Técnicas de alignment: diferentes métodos restringen distintos canales de fallo}

Human/AI feedback ajusta preferencias de salida; process supervision verifica pasos intermedios; constitutional AI genera críticas y revisiones mediante reglas; scalable oversight intenta que supervisión limitada cubra sistemas más potentes; interpretability busca mecanismos internos peligrosos. Son complementarias, no sustitutos: las reglas pueden entrar en conflicto, el feedback de IA puede heredar sesgos, las etiquetas de proceso pueden recompensar formato superficial e interpretabilidad también puede ser incompleta.

\lecturefigure{slide-148.jpg}{Técnicas de alignment I: feedback, preferencia, interpretabilidad y process supervision}{CS25 V6 oficial deck, p.~148}

\lecturefigure{slide-149.jpg}{Técnicas de alignment II: constitutional AI, scalable oversight y human feedback}{CS25 V6 oficial deck, p.~149}

\teachervoice{01:11:00--01:13:02, el docente coloca constitutional AI, scalable oversight, process supervision e interpretabilidad en la misma caja de herramientas, indicando explícitamente que el human feedback por sí mismo no garantiza alignment.}

\begin{knowledgebox}{Matriz de métodos de alignment}
\begin{tabularx}{\textwidth}{p{0.2\textwidth}p{0.25\textwidth}X}
\toprule
Método & Objeto principal de restricción & Problemas sin resolver \\
\midrule
Preference feedback & Clasificación y estilo de salida & Representatividad de preferencias, reward hacking \\
Process supervision & Pasos intermedios & Múltiples rutas equivalentes, costo de etiquetado, adaptación superficial \\
Constitutional AI & Principios explícitos & Conflictos de reglas, brecha entre explicación y ejecución \\
Scalable oversight & Supervisión de sistemas potentes & Cómo un supervisor débil juzga fiablemente tareas complejas \\
Interpretability & Feature/circuit internos & Cobertura incompleta, explicaciones engañosas \\
\bottomrule
\end{tabularx}
\end{knowledgebox}

La evaluación de alignment también requiere estratificación. Capability eval mide si el sistema puede completar tareas; alignment eval verifica si sigue eligiendo comportamientos permitidos ante conflictos de objetivos, inducción, límites de permisos y nuevas distribuciones; no son sustituibles entre sí. El conjunto de pruebas debe incluir solicitudes normales, adversarial prompt, herramientas que retornan contenido malicioso, tentaciones locales en trayectorias largas y escenarios de shutdown/abstention. También se debe reportar false refusal, ya que rechazar todas las solicitudes difíciles puede mejorar ciertas métricas de seguridad pero dañar la utilidad. Finalmente, es necesario trazar el frente de Pareto helpfulness--safety--cost, no ocultar compromisos con una única "puntuación de seguridad".

\subsection{Resumen de la sección}

La gestión de memory controla estados externos; continual learning modifica adaptación a largo plazo; model editing intenta escritura local de pesos; interpretabilidad provee evidencia interna; alignment restringe objetivos y comportamientos. No son sustituibles: un sistema con base de recuperación no necesariamente aprende continuamente, y uno que explica la atención no es necesariamente seguro.

